\section{Considerations for Release}
\label{sec:release}
Following the recommendations for individual researchers generated by the Partnership for AI,\footnote{\url{https://partnershiponai.org/paper/responsible-publication-recommendations/}} along with the governance guidance outlined by NIST,\footnote{\url{https://nvlpubs.nist.gov/nistpubs/SpecialPublications/NIST.SP.1270.pdf}} we are disclosing all of the details involved in training OPT-175B through our logbook,\footnote{\url{https://github.com/facebookresearch/metaseq/blob/main/projects/OPT/chronicles/OPT175B_Logbook.pdf}} our code, and providing researchers access to model weights for \OPT{}, along with a suite of smaller baselines mirroring the setup for \OPT{}.  We aim to be fully accountable for the development lifecycle of \OPT{}, and only through increasing transparency around LLM development can we start understanding the limitations and risks of LLMs before broader deployment occurs.

By sharing a detailed account of our day-to-day training process, we disclose not only how much compute was used to train the current version of OPT-175B, but also the human overhead required when underlying infrastructure or the training process itself becomes unstable at scale.  These details are generally omitted from previous publications, likely due to the inability to fully ablate changes made mid-flight (without drastically increasing the compute budget).  We hope that by revealing how certain ad-hoc design decisions were made, we can improve upon these practices in the future, and collectively increase the experimental robustness in developing models at this scale.

Outside of these notes, the metaseq codebase itself is the final source of truth in many of our implementation details.  By releasing our development codebase, we aim to shed light on any implementation detail that may have been omitted from being explicitly enumerated in this paper, as it is either considered a detail of standard practice in the field, or is simply a detail we failed to account for.  This current codebase is also the only known open-source implementation of training a decoder-only transformer that is $\geq$175B parameters without the use of pipeline paralellism on NVIDIA GPUs.

To enable experimentation at 175B scale, we are providing researchers with direct access to the parameters of \OPT{}. The reasoning here is two-fold: enable Responsible AI research into LLMs while simultaneously reducing the environmental impact of pursuing research at this scale.  There is a growing body of work detailing ethical and social risks from deploying language models with emergent capabilities at scale \cite{deepmindLMrisks2021,foundation2021,dinan2021anticipating,alignmentagents2021}. By limiting access to \OPT{} to the research community with a non-commercial license, we aim to focus development efforts on quantifying the limitations of the LLMs first, before broader commercial deployment occurs.
 
Furthermore, there exists significant compute and carbon cost to reproduce models of this size.  While \OPT{} was developed with an estimated carbon emissions footprint (CO2eq) of 75 tons,\footnote{With ablations, baselines and downtime, our own estimates of total cost is roughly 2$\times$ higher.} \biggpt{} was estimated to use 500 tons \cite{patterson2021carbon}, while Gopher required 380 tons \cite{gopher2022}.  These estimates are not universally reported, and the accounting methodologies for these calculations are also not standardized. In addition, model training is only one component of the overall carbon footprint of AI systems; we must also consider experimentation and eventual downstream inference cost, all of which contribute to the growing energy footprint of creating large-scale models \cite{sustainable_ai2021}.  By releasing our logbook, we hope to highlight the gap between a theoretical carbon cost estimate that assumes no hardware failures or training instabilities, versus one that aims to include the entire LLM development lifecycle.  We need to understand the manufacturing (or embodied) carbon of these systems \cite{gupta2021chasing} as they grow increasingly more complex, and we hope that our paper can help future work in defining additional factors to consider when measuring the impact of scale on the environment.


Similarly, by producing a set of baselines across a wide range of scales, we hope to enable the broader research community to study the impact and limitations of these models with respect to scale alone. As reported in \citet{chinchilla2022}, many of these LLMs may have been under-trained as a function of the amount of training data used, which implies that incorporating more data and continuing to train these baseline models may continue to improve performance.  There is also evidence that step-function changes in capabilities may occur at a scale that is much smaller than 175B \cite{flan2021}, indicating a need to examine a wider range of scales for different research applications.
